\documentclass[10pt,a4paper]{amsart}
\usepackage[margin=19mm]{geometry}
\usepackage{amsmath,amssymb,mathtools}
\usepackage[scr=rsfs,cal=euler]{mathalfa}
\usepackage{xfrac}
\usepackage[unicode,hidelinks,bookmarks=false]{hyperref}
\newcommand{\ie}{i.e.\ }
\newcommand{\cf}{cf.\ }
\newcommand{\resp}{respectively\ }
\newcommand{\Subsection}{\subsection}
\providecommand{\nobreakdash}{\nobreak\hbox{-}}
\setlength{\emergencystretch}{2em}
\multlinegap0pt
\usepackage{bm}
\usepackage{bbm}
\usepackage{dsfont}
\usepackage{xspace}
\usepackage{cancel}
\usepackage{mathtools}
\usepackage{amsthm}
\makeatletter
\@ifundefined{theorem}{\newtheorem{theorem}{Theorem}}{}
\@ifundefined{lemma}{\newtheorem{lemma}[theorem]{Lemma}}{}
\@ifundefined{proposition}{\newtheorem{proposition}[theorem]{Proposition}}{}
\makeatother
\newcommand{\R}{\mathbb{R}}
\title[The non-relativistic limit of scattering states for the Vlas]{The non-relativistic limit of scattering states for the Vlasov equation with short-range interaction potentials}
\author{Younghun Hong, Stephen Pankavich}
\date{Source: arXiv:2509.08072v1 (CC BY 4.0)}
\begin{document}
\maketitle

\noindent\emph{Source and licence.} The non-relativistic limit of scattering states for the Vlasov equation with short-range interaction potentials. Version arXiv:2509.08072v1 (CC BY 4.0), distributed under the Creative Commons Attribution 4.0 International licence, as recorded in the arXiv licence metadata for this exact version and shown on the abstract page https://arxiv.org/abs/2509.08072v1. Full rights notice: licenses/arxiv-2509.08072-CC-BY-4.0.txt.\par\smallskip

\noindent\textit{Editorial scope.} Counted unit: the Lemma whose source label is \texttt{lemma: difference between force fields} in the article (source0.tex lines 1104--1109, source label \texttt{lemma: difference between force fields}), delivered together with the article's own complete original proof of it (source0.tex lines 1131--1275, one \texttt{proof} environment of 145 lines). No mathematics is written, repaired or re-derived: statement and proof are the article's own text line for line. The proof is detached from the statement in the source: the article states the result at lines 1104--1109 and prints its proof at lines 1131--1275, and the proof environment's own header names the statement, so the association is the article's own. Required context retained uncounted: eq: rVP (lines 120--127); Adef (lines 164--167); T1 (lines 208--219); eq: initial data assumption (lines 211--213); eq: uniform force field decay bounds (lines 215--217); T2 (lines 234--244); eq: Ac asymptotic (lines 344--349); lemma: Ac eigenvalues (lines 373--376); lemma: dispersive bounds for the free flow associated with interpolated v (lines 400--410); lemma: interpolation inequality (lines 439--443); eq: small force field assumption (lines 458--460); eq: r-ODE (lines 464--471); lemma: perturbed momentum (lines 490--508); eq: Ac difference bound (lines 499--502); lemma:  the wave operator, almost identity (lines 585--597); eq: approx identity, inverse finite-time wave operator (lines 594--596); prop: nr limit for Vlasov (lines 1040--1057); eq: interpolated wave operator (lines 1112--1114); lemma: interpolated wave operator (lines 1117--1122). Every definition, display and label of those lines is retained unchanged. Omitted from this package and not redistributed: 1--119, 128--163, 168--207, 214--214, 218--233, 245--343, 350--372, 377--399, 411--438, 444--457, 461--463, 472--489, 503--584, 597--1039, 1058--1103, 1110--1111, 1115--1116, 1123--1130, 1276--1440. Nothing inside a retained excerpt is omitted or altered: every retained body is delivered exactly as the article's lines are written, so no omission receipt of any kind accompanies this package. Citations: the retained excerpts carry no citation command at all, so no bibliography entry of the article is transcribed and no reference is redistributed. Reference adaptation: none was needed and none was made; every reference inside the delivered unit resolves inside it, so no unresolved reference or citation is produced. Printed slips kept literal: no printed word, symbol, hypothesis or number of the counted statement or of its proof was altered. Layout and class: the article's own class is \texttt{amsart} with packages including amsmath, amsfonts, amssymb, graphicx, amsthm, color, yfonts, cite, paralist, hyperref, physics, dsfont, soul, wrapfig; the wrapper is the lane's pinned \texttt{amsart} editorial wrapper at 10pt on A4, which restates the theorem-like environments the retained text needs and adds the article's own macro definitions that the retained text needs (\texttt{\textbackslash R}), with extra wrapper packages (bm, bbm, dsfont, xspace, cancel, mathtools). The delivered numbering is the unit's own. Assets: no retained span contains a figure environment, a graphics-inclusion command, a PDF-page inclusion, a table or a TikZ picture; no image file is copied into this package, the package declares no asset dependency and no image file is redistributed with it. Licence: version arXiv:2509.08072v1 of this article is distributed under the Creative Commons Attribution 4.0 International licence, as recorded in the arXiv licence metadata for this exact version and shown on the abstract page https://arxiv.org/abs/2509.08072v1. A full rights notice is delivered at licenses/arxiv-2509.08072-CC-BY-4.0.txt. Characters: every retained excerpt is pure ASCII, so the delivered engine, XeLaTeX with its default Latin Modern text font, reports no missing character.
\begin{equation}\label{eq: rVP}
\left.
\begin{aligned}
\partial_t f_c+v_c(p)\cdot\nabla_x f_c+E_c\cdot\nabla_p f_c&=0,\\
f_c(0)&=f^0,
\end{aligned}
\right\}
\end{equation}

\begin{equation}
\label{Adef}
    \mathbb{A}_c(p):=\nabla v_c(p)=\frac{1}{\gamma_c(p)}\mathbb{I}_3-\frac{1}{\gamma_c(p)^3}\bigg(\frac{p_i p_j}{c^2}\bigg)_{i,j=1}^3,
\end{equation}

\begin{theorem}[Small data solutions]
\label{T1}
Assume
\begin{equation}\label{eq: initial data assumption}
\eta:=\Big\|\langle x\rangle^{3^+}\langle p \rangle^8 f^0\Big\|_{L_{x,p}^\infty} +\Big\|\langle x\rangle^{3^+}\langle p \rangle^9\nabla_{(x,p)}f^0\Big\|_{L_{x,p}^\infty}
\end{equation}
is sufficiently small. Then, for any $1 \leq c \leq \infty$, there exists a unique, global solution $f_c$ satisfying \eqref{eq: rVP} for all $t \in [0,\infty)$ and $(x,p) \in \mathbb{R}^6$. Moreover, the associated force field satisfies the uniform decay bounds
\begin{equation}\label{eq: uniform force field decay bounds}
\sup_{t\geq0}\Big\{(1+t)^{\alpha+1}\|E_c(t)\|_{L_x^\infty(\mathbb{R}^3)}+(1+t)^{\alpha+2}\|\nabla_x E_c(t)\|_{L_x^\infty(\mathbb{R}^3)}\Big\}\lesssim\eta_0,
\end{equation}
for some $\eta_0$ satisfying $0 < \eta < \eta_0 \ll 1$ where the implicit constant in \eqref{eq: uniform force field decay bounds} does not depend on $c\in [1,\infty]$.
\end{theorem}

\begin{theorem}[Scattering]
\label{T2}
For any $1 \leq c \leq \infty$, let $f_c$ be the unique, global solution constructed within Theorem \ref{T1}.
Then, for any $1 \leq c \leq \infty$, there exists $(\mathcal{X}_c^+,  \mathcal{P}_c^+) \in C(\mathbb{R}^6)$ and non-negative $f_c^+ \in L^1(\mathbb{R}^6)$ such that 
$$\Big(\mathcal{X}_c(t,0,x,p)-tv_c\big(\mathcal{P}_c(t,0,x,p)\big), \mathcal{P}_c(t,0,x,p)\Big) \to \big(\mathcal{X}_c^+(x,p), \mathcal{P}_c^+(x,p)\big)$$
and
$$f_c\Big(t,x+tv_c(p),p\Big) \to f_c^+(x,p)$$
as $t \to \infty$ for all $(x,p) \in \R^6$  with the convergence estimate
$$ \Big\Vert f_c\Big(t,x+tv_c(p),p\Big) - f_c^+(x,p)\Big\Vert_{L^1_{x,p}(\mathbb{R}^6)} \lesssim \frac{1}{(1+t)^\alpha}\|\nabla_{(x,p)}f^0\|_{L_{x,p}^1(\mathbb{R}^6)}$$
for all $t \geq 0$.
\end{theorem}

\begin{align}
|v_c(p)-p|&\leq|p|\min\bigg\{1,\frac{|p|^2}{c^2}\bigg\},\label{eq: vc asymptotic}\\
\bigg\|\mathbb{A}_c(p)-\frac{1}{\gamma_c(p)}\mathbb{I}_3\bigg\|&\leq\frac{1}{\gamma_c(p)}\min\bigg\{1,\frac{|p|^2}{c^2}\bigg\},\label{eq: Ac asymptotic}\\
\bigg\|\mathbb{A}_c(p)-\mathbb{I}_3\bigg\|&\lesssim\min\bigg\{1,\frac{|p|^2}{c^2}\bigg\},\label{eq: Ac asymptotic'}\\
|\partial_{p_k}\mathbb{A}_c^{ij}(p)|&\lesssim\frac{1}{c\gamma_c(p)^2}\min\bigg\{1,\frac{|p|}{c}\bigg\},\label{eq: Ac  derivative asymptotic}
\end{align}

\begin{lemma}\label{lemma: Ac eigenvalues}
The matrix $\nabla v_c^\theta(p)$ has 3 eigenvalues: $\frac{\theta}{\gamma_c(p)}+(1-\theta)$ of multiplicity $2$ and $\frac{\theta}{\gamma_c(p)^3}+(1-\theta)$. Thus, $\nabla v_c^\theta(p)$ is symmetric positive-definite with
$$\textup{det}(\nabla v_c^\theta(p))=\bigg(\frac{\theta}{\gamma_c(p)}+(1-\theta)\bigg)^2\bigg(\frac{\theta}{\gamma_c(p)^3}+(1-\theta)\bigg).$$ 
\end{lemma}

\begin{lemma}[Dispersive bounds for the free flow associated with $v_c^\theta(p)$]\label{lemma: dispersive bounds for the free flow associated with interpolated v}
For all $t\geq0$, we have
\begin{equation}\label{eq: dispersive bounds for the free flow associated with interpolated v-1}
\big\|T^\theta[h](t,x)\big\|_{L_x^1(\mathbb{R}^3)}\leq \|h(t,x,p)\|_{L_{x,p}^1(\mathbb{R}^6)}
\end{equation}
and
\begin{equation}
\label{eq: dispersive bounds for the free flow associated with interpolated v-2}
\big\|T^\theta[h](t,x)\big\|_{L_x^\infty(\mathbb{R}^3)}\lesssim \frac{1}{(1+t)^3}\Big\|\langle x\rangle^{3^+}\langle p\rangle^5 h(t,x,p)\Big\|_{L_{x,p}^\infty(\mathbb{R}^6)}.
\end{equation}
\end{lemma}

\begin{lemma}[Interpolation inequality]
\label{lemma: interpolation inequality}
Assume $h \in L^1(\mathbb{R}^3) \cap L^\infty(\mathbb{R}^3)$. If $|\nabla w(x)|\lesssim\frac{1}{|x|^{\alpha+1}}$ for $0<\alpha<2$, then 
$$\|\nabla w*h\|_{L^\infty(\mathbb{R}^3)}\lesssim \|h\|_{L^1(\mathbb{R}^3)}^{\frac{2-\alpha}{3}}\|h\|_{L^\infty(\mathbb{R}^3)}^{\frac{\alpha+1}{3}}.$$
\end{lemma}

\begin{equation}\label{eq: small force field assumption}
\sup_{t\geq0}\Big((1+t)^{\alpha+1}\|E(t)\|_{L_x^\infty(\mathbb{R}^3)}+(1+t)^{\alpha+2}\|\nabla _xE(t)\|_{L_x^\infty(\mathbb{R}^3)}\Big)\leq\eta_0.
\end{equation}

\begin{equation}\label{eq: r-ODE}
\left.
\begin{aligned}
\partial_s\Xi_c(s,t,x,p)&=\big(v_c(\mathcal{P}_c(s,t,x,p)), E(s,\mathcal{X}_c(s,t,x,p))\big),\\
\Xi_c(t,t,x,p)&=(x,p)\in\mathbb{R}^3\times\mathbb{R}^3,
\end{aligned}
\right \}
\end{equation}

\begin{lemma}[Perturbed momentum]
\label{lemma: perturbed momentum}
Assume that the force field $E(t,x)$ satisfies the decay bound \eqref{eq: small force field assumption} with $\alpha > 1$ and $0< \eta_0 \ll 1$ sufficiently small. If the backward characteristic flow $\Xi_c$ solves \eqref{eq: r-ODE}, then for all $0\leq s\leq t$,
we have
\begin{equation}\label{eq: bound for Pc}
\left | \mathcal{P}_c(s,t,x,p)- p \right | \lesssim \eta_0,
\end{equation}
where the implicit constant is independent of $c \geq 1$ and $s,t \geq 0$.
Moreover, if $|p-p'|\lesssim\eta_0$, then
\begin{align}
\left |\gamma_c(p') - \gamma_c(p) \right | & \lesssim \frac{\eta_0}{c},\label{eq: gamma c difference bound}\\
\left \Vert \mathbb{A}_c(p') - \mathbb{A}_c(p) \right \Vert & \lesssim \frac{\eta_0}{c}\gamma_c(p)^{-2}.\label{eq: Ac difference bound}
\end{align}
Finally, combining equations \eqref{eq: Ac asymptotic} and \eqref{eq: Ac difference bound} yields
\begin{equation}\label{A bound}
\Vert \mathbb{A}_c(\mathcal{P}_c(s)) \Vert \lesssim \gamma_c(p)^{-1}
\end{equation}
for any $ s \geq 0.$
\end{lemma}

\begin{lemma}[Almost identity]
\label{lemma:  the wave operator, almost identity}
If $E(t,x)$ satisfies \eqref{eq: small force field assumption} , then
\begin{equation}\label{eq: approx identity, finite-time wave operator}
\sup_{t\geq 0}\Big\|\mathcal{W}_c(t)(x,p)-(x,p)\Big\|_{C_{x,p}(\mathbb{R}^6)}+\sup_{t\geq 0}\Big\|\Big(\nabla_{(x,p)}\mathcal{W}_c(t)\Big)(x,p)-\mathbb{I}_6\Big\|_{C_{x,p}(\mathbb{R}^6)}\lesssim\eta_0,
\end{equation}
where the implicit constant is independent of $1 \leq c \leq \infty$. Thus, if $\eta_0>0$ is sufficiently small, then $\mathcal{W}_c(t)$ is invertible with 
$$\mathcal{W}_c(t)^{-1} = \Phi_c(t)^{-1} \circ \Phi_c^{\textup{free}}(t)$$
and
\begin{equation}\label{eq: approx identity, inverse finite-time wave operator}
\sup_{t\geq 0}\Big\|\mathcal{W}_c(t)^{-1}(x,p)-(x,p)\Big\|_{C_{x,p}(\mathbb{R}^6)}+\sup_{t\geq 0}\Big\|\Big(\nabla_{(x,p)}\mathcal{W}_c(t)^{-1}\Big)(x,p)-\mathbb{I}_6\Big\|_{C_{x,p}(\mathbb{R}^6)}\lesssim\eta_0.
\end{equation}
\end{lemma} 

\begin{proposition}[Non-relativistic limit for the Vlasov equation]\label{prop: nr limit for Vlasov}
Under the assumptions of Theorems \ref{T1} and \ref{T2}, let
$$f_c(t,x,p)=f^0\Big(\Phi_c(t)^{-1}(x,p)\Big)$$
for every $1 \leq c \leq \infty$
be the global solution to the relativistic (or non-relativistic) Vlasov equation with initial data,  constructed in Theorem \ref{T1}. Then, for $t\geq 1$, we have
\begin{equation}\label{eq: nr limit for Vlasov}
\|g_c(t,x,p)-g_\infty(t,x,p)\|_{L_{x,p}^1(\mathbb{R}^6)}\lesssim \frac{1}{c^2} \|\langle p \rangle^3\nabla_{(x,p)}f^0\|_{L_{x,p}^1(\mathbb{R}^6)},
\end{equation}
where $g_c(t,x,p)=f_c(t,x+tv_c(p),p)$. Moreover, the wave operator and force field obey the bounds
\begin{equation}\label{eq: nr limit for wave operators}
\sup_{t\geq0}\bigg\|\frac{1}{\langle p\rangle^3}\Big(\mathcal{W}_c(t)(x,p)-\mathcal{W}_\infty(t)(x,p)\Big)\bigg\|_{C_{x,p}(\mathbb{R}^6)}\lesssim \frac{\eta_0}{c^2}
\end{equation}
and
\begin{equation}\label{eq: nr limit for force field}
\|(E_c-E_\infty)(t)\|_{L_x^\infty(\mathbb{R}^3)}\lesssim\frac{\eta_0}{c^2\langle t\rangle^{\alpha+1}},
\end{equation}
respectively.
\end{proposition}

\begin{equation}\label{eq: interpolated wave operator}
\widetilde{\mathcal{W}}^\theta(t)=\big(\widetilde{\mathcal{W}}_1^\theta(t),\widetilde{\mathcal{W}}_2^\theta(t)\big):=\theta\mathcal{W}_c(t)+(1-\theta)\mathcal{W}_\infty(t):\mathbb{R}^3\times\mathbb{R}^3\to\mathbb{R}^3\times\mathbb{R}^3
\end{equation}

\begin{lemma}[Interpolated wave operator]\label{lemma: interpolated wave operator}
Under the assumptions of Proposition \ref{prop: nr limit for Vlasov}, we have
\begin{equation}\label{eq: interpolated wave operator, almost identity}
\sup_{t\geq0}\big\|\widetilde{\mathcal{W}}^\theta(t)(x,p)-(x,p)\big\|_{C_{x,p}(\mathbb{R}^6)}+\sup_{t\geq0}\big\|\nabla_{(x,p)}\widetilde{\mathcal{W}}^\theta(t)(x,p)-\mathbb{I}_6\big\|_{C_{x,p}(\mathbb{R}^6)}\lesssim\eta_0, 
\end{equation}
\end{lemma}

\begin{lemma}\label{lemma: difference between force fields}
Under the assumptions of Proposition \ref{prop: nr limit for Vlasov}, we have
$$\begin{aligned}
\|(E_c-E_\infty)(t)\|_{L_x^\infty(\mathbb{R}^3)}&\lesssim\frac{\eta_0}{(1+t)^{\alpha+1}}\bigg\{\frac{1}{c^2}+\Big\|\frac{1}{\langle p\rangle^3}\Big(\mathcal{W}_c(t)(x,p)-\mathcal{W}_\infty(t)(x,p)\Big)\Big\|_{C_{x,p}(\mathbb{R}^6)}\bigg\}.
\end{aligned}$$
\end{lemma}

\begin{proof}[Proof of Lemma \ref{lemma: difference between force fields}]
\textbf{Step 1 - Density Estimate}

\noindent Recalling that $f_c(t,x,p)=g_c(t,x-tv_c(p),p)$ and 
$g_c(t,x,p)=f^0(\mathcal{W}_c(t)^{-1}(x,p))$ for $1\leq c\leq\infty$, we write the difference between the distribution functions as 
$$\begin{aligned}
(f_c-f_\infty)(t,x,p)&=g_c\Big(t,x-tv_c(p),p\Big)-g_\infty\Big(t,x-tp,p\Big)\\
&=(g_c-g_\infty)\Big(t,x-tv_c(p),p\Big)+\Big\{g_\infty\Big(t,x-tv_c(p),p\Big)-g_\infty\Big(t,x-tp,p\Big)\Big\}\\
&=\textup{(A)}+\textup{(B)}.
\end{aligned}$$
Then, introducing $v_c^\theta(p)=\theta v_c(p)+(1-\theta)p$, we have 
$$\textup{(B)}=\int_0^1 \frac{d}{d\theta}g_\infty\Big(t,x-tv_c^\theta(p),p\Big)d\theta=t\int_0^1 \nabla_xg_\infty\Big(t,x-tv_c^\theta(p),p\Big)\cdot\big(p-v_c(p)\big)d\theta.$$
Notice that
%\footnote{Precisely, $$\begin{aligned}
%\begin{bmatrix}
%\nabla_{p_1}\{g_\infty(\cdots)\}\\
%\nabla_{p_2}\{g_\infty(\cdots)\}\\
%\nabla_{p_3}\{g_\infty(\cdots)\}
%\end{bmatrix}&=-t\begin{bmatrix}
%\nabla_{p_1}v_{\theta;1}(p)& \nabla_{p_1}v_{\theta;2}(p)& \nabla_{p_1}v_{\theta;3}(p)\\
%\nabla_{p_2}v_{\theta;1}(p)& \nabla_{p_2}v_{\theta;2}(p)& \nabla_{p_2}v_{\theta;3}(p)\\
%\nabla_{p_3}v_{\theta;1}(p)& \nabla_{p_3}v_{\theta;2}(p)& \nabla_{p_3}v_{\theta;3}(p) \end{bmatrix}\begin{bmatrix}
%(\nabla_{x_1}g_\infty)(\cdots)\\
%(\nabla_{x_2}g_\infty)(\cdots)\\
%(\nabla_{x_3}g_\infty)(\cdots)
%\end{bmatrix}+\begin{bmatrix}
%(\nabla_{p_1}g_\infty)(\cdots)\\
%(\nabla_{p_2}g_\infty)(\cdots)\\
%(\nabla_{p_3}g_\infty)(\cdots)
%\end{bmatrix},
%\end{aligned}$$
%where $(\cdots)=(t,x-tv_c^\theta(p),p)$.}
$$\begin{aligned}
\nabla_p\Big(g_\infty\Big(t,x-tv_c^\theta(p),p\Big)\Big)&=-t\Big[\nabla_p v_c^\theta(p)\Big]\nabla_xg_\infty\Big(t,x-tv_c^\theta(p),p\Big)\\
&\quad+\nabla_pg_\infty\Big(t,x-tv_c^\theta(p),p\Big),
\end{aligned}$$
and $\nabla v_c^\theta(p)=\theta \nabla v_c(p)+(1-\theta)\mathbb{I}_3$ is invertible, because Lemma \ref{lemma: Ac eigenvalues} ensures that the symmetric matrix $\nabla v_c(p)=\mathbb{A}_c(p)$ is positive definite. Thus, we have
$$\begin{aligned}
&\nabla_xg_\infty\Big(t,x-tv_c^\theta(p),p\Big)\cdot\big(p-v_c(p)\big)\\
&=\bigg\langle\frac{1}{t}\Big[\nabla_p v_c^\theta(p)\Big]^{-1}\bigg[\nabla_p g_\infty\Big(t,x-tv_c^\theta(p),p\Big)-\nabla_p\Big(g_\infty\Big(t,x-tv_c^\theta(p),p\Big)\Big)\bigg],\big(p-v_c(p)\big)\bigg\rangle_{\mathbb{R}^3}\\
&=\frac{1}{t}\bigg\langle\bigg[\nabla_pg_\infty\Big(t,x-tv_c^\theta(p),p\Big)-\nabla_p\Big(g_\infty\Big(t,x-tv_c^\theta(p),p\Big)\Big)\bigg],\Big[\nabla_p v_c^\theta(p)\Big]^{-1}\big(p-v_c(p)\big)\bigg\rangle_{\mathbb{R}^3},
\end{aligned}$$
where both $a\cdot b$ and $\langle a, b\rangle_{\mathbb{R}^3}$ are used to denote the inner product in $\mathbb{R}^3$. We insert this within $\textup{(B)}$ and integrate $(f_c-f_\infty)(t,x,p)$ in $p$ over $\mathbb{R}^3$ while maintaining the previous expression for $(\textup{A})$. Then, upon integrating by parts in the last term, the difference between the density functions can be decomposed as 
$$\begin{aligned}
(\rho_{f_c}-\rho_{f_\infty})(t,x)&=\int_{\mathbb{R}^3}(g_c-g_\infty)\Big(t,x-tv_c(p),p\Big)dp\\
&\quad+\int_0^1\int_{\mathbb{R}^3}\nabla_pg_\infty\Big(t,x-tv_c^\theta(p),p\Big)\cdot w_\theta(p)dpd\theta\\
&\quad+\int_0^1\int_{\mathbb{R}^3}g_\infty\Big(t,x-tv_c^\theta(p),p\Big)\nabla_p\cdot w_\theta(p)dpd\theta\\
&=\textup{(I)}+\textup{(II)}+\textup{(III)},
\end{aligned}$$
where 
$$w_\theta(p):=\Big[\nabla_p v_c^\theta(p)\Big]^{-1}\big(p-v_c(p)\big)$$
For $\textup{(II)}$ and $\textup{(III)}$, we note that  
$$w_\theta(p)=\Big[\nabla_p v_c^\theta(p)\Big]^{-1}\bigg(1-\frac{1}{\gamma_c(p)}\bigg)p=\frac{1-\frac{1}{\gamma_c(p)}}{\frac{\theta}{\gamma_c(p)^3}+1-\theta}p,$$
because by the identity 
$$\mathbb{A}_c(p)p=\frac{1}{\gamma_c(p)}p-\frac{\frac{|p|^2}{c^2}}{\gamma_c(p)^3}p=\frac{1}{\gamma_c(p)^3}p,$$
which follows from \eqref{Adef}, we have
$$[\nabla_p v_c^\theta(p)]p=\theta \mathbb{A}_c(p)p+(1-\theta)p=\bigg(\frac{\theta}{\gamma_c(p)^3}+(1-\theta)\bigg)p$$
and $\nabla_p v_c^\theta(p)$ is an invertible matrix. Hence, we have 
$$|w_\theta(p)|=\frac{1-\frac{1}{\gamma_c(p)}}{\frac{\theta}{\gamma_c(p)^3}+1-\theta}|p|\leq\frac{\frac{\gamma_c(p)^2-1}{\gamma_c(p)(\gamma_c(p)+1)}}{\frac{1}{\gamma_c(p)^3}}|p|\leq\frac{\gamma_c(p)|p|^3}{c^2}$$
and its divergence 
$$\nabla_p\cdot w_\theta(p)=\frac{\frac{\nabla\gamma_c(p)}{\gamma_c(p)^2}}{\frac{\theta}{\gamma_c(p)^3}+1-\theta}\cdot p-\frac{1-\frac{1}{\gamma_c(p)}}{(\frac{\theta}{\gamma_c(p)^3}+1-\theta)^2}\frac{3\theta(\nabla\gamma_c(p))}{\gamma_c(p)^4}\cdot p+\frac{3(1-\frac{1}{\gamma_c(p)})}{\frac{\theta}{\gamma_c(p)^3}+1-\theta}$$
satisfies
$$\big|\nabla_p\cdot w_\theta(p)\big|\lesssim\frac{\gamma_c(p)|p|^2}{c^2}.$$
Thus, applying the bounds for $|w_\theta(p)|$ and $|\nabla_p\cdot w_\theta(p)|$ to $\textup{(II)}$ and $\textup{(III)}$ respectively, it follows that 
$$\begin{aligned}
\big|(\rho_{f_c}-\rho_{f_\infty})(t,x)\big|&\lesssim\int_{\mathbb{R}^3}(|g_c-g_\infty|)\Big(t,x-tv_c(p),p\Big)dp\\
&\quad+\frac{1}{c^2}\int_0^1\int_{\mathbb{R}^3}\Big(\gamma_c(p)|p|^3|\nabla_pg_\infty|\Big)\Big(t,x-tv_c^\theta(p),p\Big)dpd\theta\\
&\quad+\frac{1}{c^2}\int_0^1\int_{\mathbb{R}^3}\Big(\gamma_c(p)|p|^2|g_\infty|\Big)\Big(t,x-tv_c^\theta(p),p\Big)dpd\theta.
\end{aligned}$$
As a consequence, by Lemma \ref{lemma: dispersive bounds for the free flow associated with interpolated v}, we obtain 
$$\begin{aligned}
\big\|(\rho_{f_c}-\rho_{f_\infty})(t)\big\|_{L_x^1}&\lesssim\|(g_c-g_\infty)(t)\|_{L_{x,p}^1}+\frac{1}{c^2}\big\|\gamma_c(p)|p|^3\nabla_pg_\infty(t)\big\|_{L_{x,p}^1}\\
&\quad+\frac{1}{c^2}\big\|\gamma_c(p)|p|^2g_\infty(t)\big\|_{L_{x,p}^1}
\end{aligned}$$
and
$$\begin{aligned}
\big\|(\rho_{f_c}-\rho_{f_\infty})(t)\big\|_{L_x^\infty}&\lesssim\frac{1}{(1+t)^3}\Big\|\langle x\rangle^{3^+}\langle p\rangle^{5}(g_c-g_\infty)(t)\Big\|_{L^\infty_{x,p}}\\
&\quad+\frac{1}{(1+t)^3c^2}\Big\|\langle x\rangle^{3^+}\langle p\rangle^{5}\gamma_c(p)|p|^3\nabla_pg_\infty(t)\Big\|_{L^\infty_{x,p}}\\
&\quad+\frac{1}{(1+t)^3c^2}\Big\|\langle x\rangle^{3^+}\langle p\rangle^{5}\gamma_c(p)|p|^2g_\infty(t)\Big\|_{L^\infty_{x,p}}.
\end{aligned}$$
Using Lemma \ref{lemma:  the wave operator, almost identity} with the change of variables $(x,p)=\mathcal{W}_\infty(t)(\tilde{x},\tilde{p})$ and Lemma \ref{lemma: perturbed momentum}, we find
$$\begin{aligned}
\big\|\gamma_c(p)|p|^3\nabla_pg_\infty(t)\big\|_{L_{x,p}^1}&=\Big\|\gamma_c(p)|p|^3(\nabla_{(x,p)}f^0)\Big(\mathcal{W}_\infty(t)^{-1}(x,p)\Big)\cdot\nabla_p\big(\mathcal{W}_\infty(t)^{-1}\big)(x,p)\Big\|_{L_{x,p}^1}\\
&\lesssim \Big\|\gamma_c\big(p(\tilde{x},\tilde{p})\big)|p(\tilde{x},\tilde{p})|^3(\nabla_{(x,p)}f^0)(\tilde{x},\tilde{p})\Big\|_{L_{\tilde{x},\tilde{p}}^1}\\
&\lesssim \big\|\gamma_c(p)|p|^3\nabla_{(x,p)}f^0\big\|_{L_{x,p}^1}.
\end{aligned}$$
In the same manner, one can show
$$\begin{aligned}
\big\|\gamma_c(p)|p|^2g_\infty(t)\big\|_{L_{x,p}^1}&\lesssim \big\|\gamma_c(p)|p|^2f^0\big\|_{L_{x,p}^1}\leq \Big\|\langle x\rangle^{3^+}\langle p\rangle^{9}\nabla_{(x,p)}f^0\Big\|_{L^\infty_{x,p}},\\
\Big\|\langle x\rangle^{3^+}\langle p\rangle^{5}\gamma_c(p)|p|^3\nabla_pg_\infty(t)\Big\|_{L^\infty_{x,p}}&\lesssim \Big\|\langle x\rangle^{3^+}\langle p\rangle^{5}\gamma_c(p)|p|^3\nabla_{(x,p)}f^0\Big\|_{L^\infty_{x,p}}
\leq \Big\|\langle x\rangle^{3^+}\langle p\rangle^{9}\nabla_{(x,p)}f^0\Big\|_{L^\infty_{x,p}},\\
\Big\|\langle x\rangle^{3^+}\langle p\rangle^{5}\gamma_c(p)|p|^2g_\infty(t)\Big\|_{L^\infty_{x,p}}&\lesssim \Big\|\langle x\rangle^{3^+}\langle p\rangle^{5}\gamma_c(p)|p|^2f^0\Big\|_{L^\infty_{x,p}}
\leq \Big\|\langle x\rangle^{3^+}\langle p\rangle^{8}f^0\Big\|_{L^\infty_{x,p}}.
\end{aligned}$$
Then, by the smallness assumption on the initial data \eqref{eq: initial data assumption}, it follows that  
\begin{equation}\label{eq: L1 Linfty rho difference bounds}
\begin{aligned}
\big\|(\rho_{f_c}-\rho_{f_\infty})(t)\big\|_{L_x^1}&\lesssim\|g_c-g_\infty\|_{L_{x,p}^1}+\frac{\eta_0}{c^2},\\
\big\|(\rho_{f_c}-\rho_{f_\infty})(t)\big\|_{L_x^\infty}&\lesssim\frac{1}{(1+t)^3}\bigg(\Big\|\Big(\langle x\rangle^{3^+}\langle p\rangle^5(g_c-g_\infty)(t)\Big\|_{L^\infty_{x,p}}+\frac{\eta_0}{c^2}\bigg).
\end{aligned}
\end{equation}
\textbf{Step 2 - Difference of Distributions} 

\noindent Next, we estimate $\|g_c-g_\infty\|_{L_{x,p}^1}$ and
$\|\langle x\rangle^{3^+}\langle p\rangle^{5}(g_c-g_\infty)(t)\|_{L^\infty_{x,p}}$
in the upper bounds \eqref{eq: L1 Linfty rho difference bounds} remaining from the previous step. To do so, we first note that 
$$g_c(t,x,p)=g_\infty\Big(t,\big(\mathcal{W}_\infty(t)\circ\mathcal{W}_c(t)^{-1}\big)(x,p)\Big).$$
Hence, using the volume-preserving property of $\mathcal{W}_c(t)$ with the interpolated wave operator \eqref{eq: interpolated wave operator}, it follows that 
$$\begin{aligned}
\|(g_c-g_\infty)
(t)\|_{L_{x,p}^1}&=\bigg\|\int_0^1 \frac{d}{d\theta}\Big[g_\infty\Big (t, \widetilde{\mathcal{W}}^\theta(t)(x,p)\Big)\Big] d\theta\bigg\|_{L_{x,p}^1}\\
&\leq\int_0^1\Big\|\nabla_{(x,p)}g_\infty\Big(t,\widetilde{\mathcal{W}}^\theta(t)(x,p)\Big)\cdot\Big(\mathcal{W}_c(t)-\mathcal{W}_\infty(t)\Big)(x,p)\Big\|_{L_{x,p}^1}d\theta\\
&\leq\bigg\|\frac{\mathcal{W}_c(t)-\mathcal{W}_\infty(t)}{\langle p\rangle^3}\bigg\|_{C_{x,p}}\int_0^1\Big\|\langle p\rangle^3\nabla_{(x,p)}g_\infty\Big(t,\widetilde{\mathcal{W}}^\theta(t)(x,p)\Big)\Big\|_{L_{x,p}^1}d\theta.
\end{aligned}$$
For the second factor in the upper bound, we use Lemma \ref{lemma: interpolated wave operator} and change variables via $$(\tilde{x},\tilde{p})=\widetilde{\mathcal{W}}^\theta(t)(x,p)=\Big(\widetilde{\mathcal{W}}_1^\theta(t)(x,p), \widetilde{\mathcal{W}}_2^\theta(t)(x,p)\Big): \mathbb{R}^3\times\mathbb{R}^3\to \mathbb{R}^3\times\mathbb{R}^3,$$ 
to find 
$$\Big\|\langle p\rangle^3\nabla_{(x,p)}g_\infty\Big(t,\widetilde{\mathcal{W}}^\theta(t)(x,p)\Big)\Big\|_{L_{x,p}^1}\lesssim\big\|\langle \tilde{p}\rangle^3\nabla_{(x,p)}g_\infty(t, \tilde{x}, \tilde{p})\big\|_{L_{\tilde{x},\tilde{p}}^1}.$$
Subsequently, applying the chain rule to $g_\infty(t,\tilde{x},\tilde{p})=f^0(\mathcal{W}_\infty(t)^{-1}(\tilde{x},\tilde{p}))$ and using \eqref{eq: approx identity, inverse finite-time wave operator}, we obtain
$$\begin{aligned}
&\Big\|\langle p\rangle^3\nabla_{(x,p)}g_\infty\Big(t,\widetilde{\mathcal{W}}^\theta(t)(x,p)\Big)\Big\|_{L_{x,p}^1}\\
&\lesssim\Big\|\langle \tilde{p}\rangle^3\nabla_{(x,p)}f^0\Big(\mathcal{W}_\infty(t)^{-1}(\tilde{x},\tilde{p})\Big)\Big\|_{L_{\tilde{x},\tilde{p}}^1}\Big\|\nabla_{(\tilde{x},\tilde{p})}\Big(\mathcal{W}_\infty(t)^{-1}(\tilde{x},\tilde{p})\Big)\Big\|_{C_{\tilde{x},\tilde{p}}}\\
&\lesssim \big\|\langle p\rangle^3\nabla_{(x,p)}f^0 \big\|_{L_{x,p}^1}
\leq \Big\|\langle x\rangle^{3^+}\langle p \rangle^9\nabla_{(x,p)}f^0\Big\|_{L_{x,p}^\infty} \leq\eta_0.
\end{aligned}$$
Therefore, we find
\begin{equation}\label{eq: difference between force fields-1}
\|(g_c-g_\infty)(t)\|_{L_{x,p}^1}\lesssim \eta_0\Big\|\frac{1}{\langle p\rangle^3}\Big(\mathcal{W}_c(t)(x,p)-\mathcal{W}_\infty(t)(x,p)\Big)\Big\|_{C_{x,p}}.
\end{equation}
Using the same tools, it follows that 
\begin{equation}\label{eq: difference between force fields-2}
\Big\|\langle x\rangle^{3^+}\langle p\rangle^5(g_c-g_\infty)(t)\Big\|_{L^\infty_{x,p}}\lesssim \eta_0\Big\|\frac{1}{\langle p\rangle^3}\Big(\mathcal{W}_c(t)(x,p)-\mathcal{W}_\infty(t)(x,p)\Big)\Big\|_{C_{x,p}}
\end{equation}
as
$\Big\|\langle x\rangle^{3^+}\langle p \rangle^8\nabla_{(x,p)}f^0\Big\|_{L_{x,p}^\infty}\leq \eta_0.$
Finally, inserting \eqref{eq: difference between force fields-1} and \eqref{eq: difference between force fields-2} within \eqref{eq: L1 Linfty rho difference bounds}, we arrive at the estimates
$$
\begin{aligned}
\big\|(\rho_{f_c}-\rho_{f_\infty})(t)\big\|_{L_x^1}&\lesssim
\eta_0 \bigg(\frac{1}{c^2} + \Big\|\frac{1}{\langle p\rangle^3}\Big(\mathcal{W}_c(t)(x,p)-\mathcal{W}_\infty(t)(x,p)\Big)\Big\|_{C_{x,p}} \bigg),\\
\big\|(\rho_{f_c}-\rho_{f_\infty})(t)\big\|_{L_x^\infty}&\lesssim
\frac{\eta_0}{(1+t)^3}\bigg(\frac{1}{c^2} + \Big\|\frac{1}{\langle p\rangle^3}\Big(\mathcal{W}_c(t)(x,p)-\mathcal{W}_\infty(t)(x,p)\Big)\Big\|_{C_{x,p}} \bigg).
\end{aligned}
$$
Using the interpolation inequality in Lemma \ref{lemma: interpolation inequality} for the difference $(\rho_{f_c}-\rho_{f_\infty})(t)$, the proof is complete.
\end{proof}

\end{document}
